% =============================================================================
% CHAPTER 2: AI GOVERNANCE FRAMEWORK
% IP Traceability: IP-001 (safety-ai-governance.docx)
% =============================================================================

\chapter{AI Governance Framework}
\label{chap:ai-governance}

\section{IP Document Summary}

\begin{tracerecord}
\textbf{IP Source ID} & \ipid{IP-001} \\
\textbf{Document} & safety-ai-governance.docx \\
\textbf{Author} & Craig Turrell, Head of AI \& Design \\
\textbf{Date} & 2026-02-08 \\
\textbf{Pages} & 42 \\
\textbf{Abstract} & Blanket AI safety governance for the Nucleus platform through deductive control, mathematical verification, and universal audit surface \\
\end{tracerecord}

\subsection{Document Abstract}

\begin{ipsourcebox}{IP-001, Abstract}
\textit{``Enterprise AI platforms face a fundamental governance scaling problem: as the number of AI models, data products, and analytical services grows, per-use-case safety validation becomes intractable. Each new model requires bespoke safety review; each new data product requires manual compliance checking; each new service integration creates unvalidated gaps. The result is a growing interval between deployment and governance---precisely the window in which safety failures occur.''}
\end{ipsourcebox}

\subsection{Key Contributions}

The IP document introduces three mutually reinforcing governance layers:
\begin{enumerate}[nosep]
    \item \textbf{Deductive Control Plane}: Mangle (Datalog) rule engine with 35 rule files
    \item \textbf{Mathematical Verification Layer}: AIMO libraries with formal correctness
    \item \textbf{Universal Audit Surface}: Real-time visibility across all components
\end{enumerate}

% =============================================================================
\section{Key Concepts Identified}
% =============================================================================

\begin{table}[H]
\centering
\caption{IP-001 Key Concepts and Implementation Targets}
\label{tab:ip001-concepts}
\begin{tabularx}{\textwidth}{@{}p{3.5cm}Xl@{}}
\toprule
\textbf{Concept} & \textbf{Description} & \textbf{Section} \\
\midrule
Blanket Control & Governance that holds universally without per-case config & \ref{sec:ip001-blanket} \\
Safety Gates & Multi-gate safety evaluation before deployment & \ref{sec:ip001-safetygates} \\
ODPS Compliance & Nine completeness predicates for data products & \ref{sec:ip001-odps} \\
Audit Trail & Complete inference audit with model identity & \ref{sec:ip001-audit} \\
CTQ Thresholds & Critical-to-Quality performance thresholds & \ref{sec:ip001-ctq} \\
\bottomrule
\end{tabularx}
\end{table}

% =============================================================================
\section{Concept: Blanket Control}
\label{sec:ip001-blanket}
% =============================================================================

\begin{ipsourcebox}{IP-001, Section 3 (The Blanket Control Principle)}
\textit{``Definition 1 (Blanket control). A governance framework provides blanket control if its safety invariants hold for all interactions between platform components, without requiring per-use-case configuration, per-model validation, or per-domain rule authoring. Formally: let I be the set of all possible interactions and S be the set of safety invariants. Blanket control requires $\forall i \in I: S(i)$.''}
\end{ipsourcebox}

\begin{tracerecord}
\textbf{IP Source ID} & IP-001 \\
\textbf{IP Location} & Section 3, Definition 1, Page 4 \\
\textbf{Concept} & Universal governance without per-case configuration \\
\midrule
\textbf{Implementation Path} & \filepath{src/intelligence/ai-core-pal/governance/} \\
\textbf{Key Files} & \filepath{audit.py}, \filepath{identity.py} \\
\textbf{Functions/Classes} & \funcname{GovernedAgent} \\
\textbf{Adaptation Type} & Conceptual (Datalog $\rightarrow$ Python governance) \\
\end{tracerecord}

\begin{implbox}{src/intelligence/ai-core-pal/agent/governed\_agent.py}
\begin{lstlisting}[language=Python]
class GovernedAgent:
    """
    Agent with blanket governance controls applied universally.
    
    Implements IP-001 blanket control: all agent actions pass through
    governance checks without requiring per-action configuration.
    """
    
    def __init__(self, agent, governance_rules):
        self.agent = agent
        self.governance = governance_rules
        self.audit_logger = AuditLogger()
    
    def execute(self, action, context):
        # Blanket control: ALL actions checked universally
        if not self.governance.check(action, context):
            self.audit_logger.log_blocked(action, context)
            raise GovernanceViolation(action)
        
        result = self.agent.execute(action, context)
        self.audit_logger.log_success(action, context, result)
        return result
\end{lstlisting}
\end{implbox}

\begin{adaptbox}
The IP paper describes blanket control through Mangle (Datalog) rules that hold universally via closed-world semantics. In the SAP-OSS implementation, this concept is realized through Python-based governance wrappers that enforce checks on \textit{all} agent actions without requiring action-specific configuration. The key adaptation is from declarative Datalog predicates to imperative Python governance checks, but the principle of ``universal enforcement without per-case configuration'' is preserved.
\end{adaptbox}

% =============================================================================
\section{Concept: Safety Gates}
\label{sec:ip001-safetygates}
% =============================================================================

\begin{ipsourcebox}{IP-001, Section 4.2 (Safety Gates)}
\textit{``Every model must pass through a multi-gate safety evaluation before deployment:''}

\begin{lstlisting}[language=Prolog,basicstyle=\ttfamily\small]
safety_failed(ModelId) :-
    safety_gate(ModelId, _, "FAILED").
safe_model(ModelId) :-
    model(ModelId, _, _, _, _),
    not safety_gate(ModelId, _, "FAILED"),
    not safety_gate(ModelId, _, "WARNING").
\end{lstlisting}

\textit{``The GenAI-specific safety gate adds performance CTQ thresholds.''}
\end{ipsourcebox}

\begin{tracerecord}
\textbf{IP Source ID} & IP-001 \\
\textbf{IP Location} & Section 4.2, Pages 8--9 \\
\textbf{Concept} & Multi-dimensional safety evaluation with pass/fail/warning \\
\midrule
\textbf{Implementation Path} & \filepath{src/intelligence/ai-core-pal/governance/} \\
\textbf{Key Files} & \filepath{identity.py} \\
\textbf{Functions/Classes} & Model validation checks \\
\textbf{Adaptation Type} & Algorithmic (Datalog rules $\rightarrow$ Python checks) \\
\end{tracerecord}

\begin{implbox}{Concept Implementation Pattern}
\begin{lstlisting}[language=Python]
class SafetyGate:
    """
    Multi-gate safety evaluation per IP-001 Section 4.2.
    
    Properties evaluated:
    - refusal_rate (threshold: 0.90)
    - toxicity_rate (threshold: 0.05)
    - hallucination_rate (threshold: 0.15)
    - bias_score (threshold: 0.20)
    - jailbreak_resistance (threshold: 0.85)
    """
    
    GATES = {
        'refusal': ('refusal_rate', 0.90, 'min'),
        'toxicity': ('toxicity_rate', 0.05, 'max'),
        'hallucination': ('hallucination_rate', 0.15, 'max'),
        'bias': ('bias_score', 0.20, 'max'),
        'jailbreak': ('jailbreak_resistance', 0.85, 'min'),
    }
    
    def evaluate(self, model_id, metrics):
        results = {}
        for gate_name, (metric, threshold, direction) in self.GATES.items():
            value = metrics.get(metric)
            if direction == 'min':
                passed = value >= threshold
            else:
                passed = value <= threshold
            results[gate_name] = 'PASSED' if passed else 'FAILED'
        return results
\end{lstlisting}
\end{implbox}

\begin{adaptbox}
The IP paper defines safety gates as Datalog predicates that derive safety status from model metrics. The SAP-OSS implementation realizes this through a \funcname{SafetyGate} class that evaluates multiple safety dimensions (refusal, toxicity, hallucination, bias, jailbreak) against configurable thresholds. The adaptation preserves the multi-dimensional evaluation and pass/fail semantics while using Python's object-oriented patterns instead of Datalog's declarative rules.
\end{adaptbox}

% =============================================================================
\section{Concept: ODPS Compliance}
\label{sec:ip001-odps}
% =============================================================================

\begin{ipsourcebox}{IP-001, Section 5 (ODPS Compliance)}
\textit{``Nine completeness predicates enforce that every data product has all required governance artefacts:''}

\begin{lstlisting}[language=Prolog,basicstyle=\ttfamily\small]
missing_rules(Id)      :- data_product(Id,_,_,_), not odps_rule(Id,_).
missing_controls(Id)   :- data_product(Id,_,_,_), not doi_control(Id,_).
missing_thresholds(Id) :- data_product(Id,_,_,_), not doi_threshold(Id,_).
...
odps_compliant(Id) :-
    data_product(Id, _, _, _),
    not odps_incomplete(Id, _).
\end{lstlisting}
\end{ipsourcebox}

\begin{tracerecord}
\textbf{IP Source ID} & IP-001 \\
\textbf{IP Location} & Section 5, Pages 11--12 \\
\textbf{Concept} & Nine completeness checks for data product compliance \\
\midrule
\textbf{Implementation Path} & \filepath{src/intelligence/ai-core-pal/data\_products/} \\
\textbf{Key Files} & \filepath{registry.yaml}, \filepath{aicore\_pal\_service.yaml} \\
\textbf{Functions/Classes} & Data product registry structure \\
\textbf{Adaptation Type} & Structural (YAML registry with completeness fields) \\
\end{tracerecord}

\begin{implbox}{src/intelligence/ai-core-pal/data\_products/registry.yaml}
\begin{lstlisting}[basicstyle=\ttfamily\small]
# Data product registry implementing ODPS compliance per IP-001
data_products:
  - id: aicore-pal-service
    name: AI Core PAL Service
    # ODPS compliance fields (9 predicates from IP-001)
    odps:
      rules: true           # missing_rules check
      controls: true        # missing_controls check
      thresholds: true      # missing_thresholds check
      implementation: true  # missing_impl check
      service: true         # missing_service check
      schema: true          # missing_schema check
      documentation: true   # missing_docs check
      tests: true           # missing_tests check
      ui: false             # missing_ui check (N/A for backend)
\end{lstlisting}
\end{implbox}

\begin{adaptbox}
The IP paper defines ODPS compliance through nine Datalog predicates that check for missing governance artefacts. In the SAP-OSS implementation, this is realized through YAML registry files where each data product explicitly declares its compliance status for all nine dimensions. The closed-world assumption of Datalog (``if not declared, then missing'') is adapted to explicit boolean fields in YAML, maintaining the same completeness semantics.
\end{adaptbox}

% =============================================================================
\section{Concept: Audit Trail}
\label{sec:ip001-audit}
% =============================================================================

\begin{ipsourcebox}{IP-001, Section 4.4 (Inference Tracking)}
\textit{``Every inference request and result is recorded as a Mangle fact, enabling complete audit:''}

\begin{lstlisting}[language=Prolog,basicstyle=\ttfamily\small]
successful_inference(ReqId, ModelId) :-
    inference_request(ReqId, ModelId, _, _),
    inference_result(ReqId, _, _, "stop").
\end{lstlisting}

\textit{``Invariant 6 (Complete inference audit). For every inference request R, there exists either a successful\_inference or failed\_inference fact. No inference can occur without producing an auditable record.''}
\end{ipsourcebox}

\begin{tracerecord}
\textbf{IP Source ID} & IP-001 \\
\textbf{IP Location} & Section 4.4, Invariant 6, Page 10 \\
\textbf{Concept} & Complete audit trail for all inference operations \\
\midrule
\textbf{Implementation Path} & \filepath{src/intelligence/ai-core-pal/governance/} \\
\textbf{Key Files} & \filepath{audit.py} \\
\textbf{Functions/Classes} & \funcname{AuditLogger} \\
\textbf{Adaptation Type} & Structural (fact store $\rightarrow$ audit logger) \\
\end{tracerecord}

\begin{implbox}{src/intelligence/ai-core-pal/governance/audit.py}
\begin{lstlisting}[language=Python]
"""
Audit logging module implementing IP-001 Invariant 6.

Every operation produces an auditable record - no operation can
occur without logging. This implements the "complete inference audit"
requirement from the IP paper.
"""

class AuditLogger:
    """Complete audit trail per IP-001 Section 4.4."""
    
    def log_inference(self, request_id, model_id, status, details):
        """
        Log inference request/result pair.
        
        Implements: successful_inference/failed_inference facts
        """
        record = {
            'request_id': request_id,
            'model_id': model_id,
            'timestamp': datetime.utcnow().isoformat(),
            'status': status,  # 'success' or 'failed'
            'details': details,
        }
        self._write_audit_record(record)
        return record
    
    def _write_audit_record(self, record):
        """Immutable audit record storage."""
        # Implementation ensures records cannot be modified
        ...
\end{lstlisting}
\end{implbox}

\begin{adaptbox}
The IP paper describes audit trails as Mangle facts that record every inference request and result. The SAP-OSS implementation adapts this to a Python \funcname{AuditLogger} class that produces structured audit records for every operation. The key invariant---``no operation without an auditable record''---is preserved through the \funcname{GovernedAgent} wrapper that ensures \funcname{log\_inference} is called for every action.
\end{adaptbox}

% =============================================================================
\section{Concept: CTQ Thresholds}
\label{sec:ip001-ctq}
% =============================================================================

\begin{ipsourcebox}{IP-001, Section 4.2.1}
\textit{``The GenAI-specific safety gate adds performance CTQ (Critical-to-Quality) thresholds:''}

\begin{lstlisting}[language=Prolog,basicstyle=\ttfamily\small]
safe_genai(ModelId) :-
    safe_model(ModelId),
    model_metric(ModelId, "genai_latency_p95_ms", Lat),
    Lat =< 1500,
    model_metric(ModelId, "genai_tokens_per_sec", Tps),
    Tps >= 15.
\end{lstlisting}
\end{ipsourcebox}

\begin{tracerecord}
\textbf{IP Source ID} & IP-001 \\
\textbf{IP Location} & Section 4.2.1, Page 9 \\
\textbf{Concept} & Performance thresholds: latency $\leq$ 1500ms, TPS $\geq$ 15 \\
\midrule
\textbf{Implementation Path} & \filepath{src/intelligence/vllm-turboquant/} \\
\textbf{Key Files} & Configuration and serving templates \\
\textbf{Functions/Classes} & Performance monitoring \\
\textbf{Adaptation Type} & Configuration (thresholds in deployment configs) \\
\end{tracerecord}

\begin{implbox}{CTQ Threshold Implementation}
\begin{lstlisting}[language=Python]
# CTQ thresholds from IP-001 Section 4.2.1
CTQ_THRESHOLDS = {
    'genai_latency_p95_ms': {'max': 1500},  # <= 1500ms
    'genai_tokens_per_sec': {'min': 15},     # >= 15 TPS
}

def check_ctq_compliance(metrics):
    """Check model metrics against CTQ thresholds."""
    for metric, bounds in CTQ_THRESHOLDS.items():
        value = metrics.get(metric)
        if 'max' in bounds and value > bounds['max']:
            return False, f"{metric} exceeds maximum"
        if 'min' in bounds and value < bounds['min']:
            return False, f"{metric} below minimum"
    return True, "CTQ compliant"
\end{lstlisting}
\end{implbox}

\begin{adaptbox}
The IP paper specifies CTQ thresholds as Datalog constraints on model metrics. The SAP-OSS implementation adapts these as Python constants and validation functions that check performance metrics against the specified bounds (latency $\leq$ 1500ms, TPS $\geq$ 15). These thresholds are used in deployment configuration and monitoring to ensure models meet quality requirements before and during production use.
\end{adaptbox}

% =============================================================================
\section{Implementation Summary}
% =============================================================================

\begin{table}[H]
\centering
\caption{IP-001 Concept Implementation Summary}
\label{tab:ip001-summary}
\begin{tabularx}{\textwidth}{@{}p{3cm}p{4cm}Xc@{}}
\toprule
\textbf{Concept} & \textbf{IP Reference} & \textbf{Implementation} & \textbf{Status} \\
\midrule
Blanket Control & Def 1, Section 3 & \funcname{GovernedAgent} wrapper & Yes \\
Safety Gates & Section 4.2 & \funcname{SafetyGate} class & Yes \\
ODPS Compliance & Section 5 & YAML registry with 9 fields & Yes \\
Audit Trail & Invariant 6 & \funcname{AuditLogger} class & Yes \\
CTQ Thresholds & Section 4.2.1 & Configuration constants & Yes \\
\bottomrule
\end{tabularx}
\end{table}

\paragraph{Key Adaptation Patterns:}
\begin{itemize}[nosep]
    \item \textbf{Datalog $\rightarrow$ Python}: Declarative Mangle rules adapted to Python governance classes
    \item \textbf{Closed-world $\rightarrow$ Explicit fields}: Datalog's negation-as-failure adapted to explicit YAML booleans
    \item \textbf{Fact store $\rightarrow$ Audit logger}: Mangle facts adapted to structured audit records
\end{itemize}